\documentclass[12pt,a4paper]{article}
\usepackage[utf8]{inputenc}
\usepackage[T1]{fontenc}
\usepackage{amsmath,amssymb}
\usepackage{enumitem}
\usepackage{geometry}
\usepackage{array}
\usepackage{multicol}
\geometry{margin=2.2cm}
\setlist{itemsep=0.25em, topsep=0.3em}
\title{Aula 10 -- SQL}
\author{Banco de Dados -- Guia de Revisão}
\date{}
\begin{document}
\maketitle
\tableofcontents
\newpage

\section{Objetivo da aula}
Esta aula apresenta SQL como linguagem padrão para definição, consulta e manipulação de dados. Para prova objetiva, o foco é identificar comandos, interpretar consultas e prever resultados.

\section{SQL como linguagem declarativa}
SQL é declarativa: o usuário diz \textbf{o que} deseja, e o SGBD decide \textbf{como} executar.

Exemplo:
\begin{verbatim}
SELECT nome
FROM Aluno
WHERE curso = 'BTI';
\end{verbatim}
A consulta diz qual resultado quer, não qual algoritmo usar.

\section{Categorias de comandos}
\begin{itemize}
    \item \textbf{DDL}: definição de dados. Ex.: \texttt{CREATE TABLE}, \texttt{ALTER TABLE}, \texttt{DROP TABLE}.
    \item \textbf{DML}: manipulação de dados. Ex.: \texttt{SELECT}, \texttt{INSERT}, \texttt{UPDATE}, \texttt{DELETE}.
    \item \textbf{VDL}: definição de visões. Ex.: \texttt{CREATE VIEW}.
\end{itemize}

\section{Vocabulário relacional em SQL}
\begin{center}
\begin{tabular}{|c|c|}
\hline
Modelo Relacional & SQL \\
\hline
Relação & Tabela \\
Tupla & Linha \\
Atributo & Coluna \\
Domínio & Tipo de dado \\
\hline
\end{tabular}
\end{center}

\section{Consulta básica}
Estrutura principal:
\begin{verbatim}
SELECT atributos
FROM tabelas
WHERE condicao;
\end{verbatim}

\begin{itemize}
    \item \texttt{SELECT}: escolhe colunas/expressões.
    \item \texttt{FROM}: escolhe tabelas.
    \item \texttt{WHERE}: filtra linhas.
\end{itemize}

\section{Exemplo visual: resultado de SELECT}
Tabela \texttt{Aluno}:
\[
\begin{array}{|c|c|c|}
\hline
\text{matricula} & \text{nome} & \text{curso}\\
\hline
1 & \text{Ana} & \text{BTI}\\
2 & \text{Bruno} & \text{Computacao}\\
3 & \text{Carla} & \text{BTI}\\
\hline
\end{array}
\]
Consulta:
\begin{verbatim}
SELECT nome
FROM Aluno
WHERE curso = 'BTI'
ORDER BY nome;
\end{verbatim}
Resultado:
\[
\begin{array}{|c|}
\hline
\text{nome}\\
\hline
\text{Ana}\\
\text{Carla}\\
\hline
\end{array}
\]

\section{DISTINCT}
Por padrão, SQL pode retornar duplicatas. Para remover repetições:
\begin{verbatim}
SELECT DISTINCT curso
FROM Aluno;
\end{verbatim}

\section{Junções}
Sem condição de junção, duas tabelas no \texttt{FROM} produzem produto cartesiano.

\subsection{Junção implícita}
\begin{verbatim}
SELECT A.nome, C.nomeCurso
FROM Aluno A, Curso C
WHERE A.idCurso = C.idCurso;
\end{verbatim}

\subsection{Junção explícita}
\begin{verbatim}
SELECT A.nome, C.nomeCurso
FROM Aluno A
JOIN Curso C ON A.idCurso = C.idCurso;
\end{verbatim}

\section{Funções agregadas}
Funções comuns:
\begin{itemize}
    \item \texttt{COUNT(*)}: conta linhas.
    \item \texttt{SUM(coluna)}: soma.
    \item \texttt{AVG(coluna)}: média.
    \item \texttt{MIN(coluna)} e \texttt{MAX(coluna)}.
\end{itemize}

\section{GROUP BY e HAVING}
\texttt{GROUP BY} agrupa linhas. \texttt{HAVING} filtra grupos.

Exemplo:
\begin{verbatim}
SELECT curso, COUNT(*)
FROM Aluno
GROUP BY curso
HAVING COUNT(*) > 5;
\end{verbatim}

Pegadinha: condição sobre agregação vai no \texttt{HAVING}, não no \texttt{WHERE}.

\section{Subconsultas}
Exemplo: alunos de cursos do centro CT.
\begin{verbatim}
SELECT nome
FROM Aluno
WHERE idCurso IN (
    SELECT idCurso
    FROM Curso
    WHERE centro = 'CT'
);
\end{verbatim}

\section{INSERT, UPDATE e DELETE}
\subsection{INSERT}
\begin{verbatim}
INSERT INTO Aluno VALUES (1, 'Ana', 'BTI');
\end{verbatim}

\subsection{UPDATE}
\begin{verbatim}
UPDATE Aluno
SET curso = 'Computacao'
WHERE matricula = 10;
\end{verbatim}
Sem \texttt{WHERE}, atualiza todas as linhas.

\subsection{DELETE}
\begin{verbatim}
DELETE FROM Aluno
WHERE matricula = 10;
\end{verbatim}
Sem \texttt{WHERE}, remove todas as linhas.

\section{Restrições de integridade}
\begin{itemize}
    \item \textbf{PRIMARY KEY}: identifica unicamente uma linha.
    \item \textbf{FOREIGN KEY}: referencia chave de outra tabela.
    \item \textbf{NOT NULL}: impede valor nulo.
    \item \textbf{UNIQUE}: impede repetição.
\end{itemize}

Exemplo:
\begin{verbatim}
CREATE TABLE Aluno (
    matricula INT PRIMARY KEY,
    nome VARCHAR(100),
    idCurso INT,
    FOREIGN KEY (idCurso) REFERENCES Curso(idCurso)
);
\end{verbatim}

\section{NULL}
Para testar \texttt{NULL}, use \texttt{IS NULL} ou \texttt{IS NOT NULL}. Não use \texttt{= NULL}.

\begin{verbatim}
SELECT nome
FROM Aluno
WHERE curso IS NULL;
\end{verbatim}

\section{Pegadinhas comuns}
\begin{itemize}
    \item \texttt{WHERE} filtra linhas; \texttt{HAVING} filtra grupos.
    \item \texttt{DELETE} remove linhas; \texttt{DROP TABLE} remove tabela.
    \item \texttt{UPDATE} sem \texttt{WHERE} altera todas as linhas.
    \item \texttt{COUNT(DISTINCT curso)} conta cursos distintos.
    \item Produto cartesiano aparece quando falta condição de junção.
\end{itemize}

\section{Desafios}
\subsection*{Desafio 1}
Considere:
\[
Aluno(\text{matricula},\text{nome},\text{idCurso})
\]
\[
Curso(\text{idCurso},\text{nomeCurso})
\]
Escreva a consulta que retorna os nomes dos alunos do curso BTI.

\subsection*{Desafio 2}
Dada a tabela \texttt{Matricula(idAluno, disciplina, nota)}, escreva uma consulta que retorna disciplinas com média maior ou igual a 7.

\section{Resolução dos desafios}
\textbf{Desafio 1}:
\begin{verbatim}
SELECT A.nome
FROM Aluno A
JOIN Curso C ON A.idCurso = C.idCurso
WHERE C.nomeCurso = 'BTI';
\end{verbatim}

\textbf{Desafio 2}:
\begin{verbatim}
SELECT disciplina, AVG(nota)
FROM Matricula
GROUP BY disciplina
HAVING AVG(nota) >= 7;
\end{verbatim}

\section{Checklist final}
\begin{itemize}
    \item Sei diferenciar DDL, DML e VDL?
    \item Sei prever resultado de \texttt{SELECT}?
    \item Sei identificar junção correta?
    \item Sei usar \texttt{GROUP BY} e \texttt{HAVING}?
    \item Sei o risco de \texttt{UPDATE} e \texttt{DELETE} sem \texttt{WHERE}?
\end{itemize}

\end{document}
